APR 05 - TEORIA CNN
CNN: rede projetada para dados com estrutura espacial/temporal local (imagens, ECG, audio).
Convolucao aplica filtros partilhados por multiplas posicoes; aprende features automaticamente.
Feature map: resultado de cada filtro. Varios filtros => varios canais de saida.
Nas primeiras camadas surgem edges, linhas e texturas; nas profundas, combinacoes e padroes abstratos.
Local receptive field: cada neuronio observa regiao local. Weight sharing reduz parametros face a dense.
Stride controla deslocamento; padding controla tamanho/bordas.
Kernel KxK normalmente aprendido por backprop.
A operacao implementada por bibliotecas costuma ser cross-correlation (sem inverter kernel).
ReLU introduz nao linearidade apos convolucao.
Max pooling resume valores mais altos; average pooling calcula media; reduzem dimensoes e sensibilidade a pequenas translacoes.
Pooling nao tem pesos treinaveis na forma comum.
Flatten transforma mapas 3D em vetor para Dense; GlobalAveragePooling resume cada canal e reduz parametros.
Conv1D para series temporais/ECG; Conv2D para imagens.
Tensor imagem (H,W,C); batch comum (N,H,W,C) no Keras.
CNN nao e invariante a toda rotacao/translacao; augmentation melhora robustez.
Classificacao ECG: CNN recebe sinal bruto normalizado, aprende padroes locais em vez de features manuais.
Comparacao MLNN: normalmente precisa de features extraidas manualmente; CNN aprende filtros.
Overfitting: usar augmentation, dropout, regularizacao, early stopping.
Output shapes e contagem de parametros sao exercicios tipicos.